% !PRINTED 487
\weberpage{487}
%<<CONT
\[
\psi_1(\alpha, y) = 0,
\]
then, on account of the transitivity of the group, we can apply to this equation a permutation by which $\alpha$ passes into an arbitrary quantity of the sequence $A$, while $y$ remains unchanged; hence every root of $\psi(t, y)$ is at the same time a root of $\psi_1(t, y)$, and therefore the two functions are identical.

Hence the equation $f(x) = 0$ of degree $n$ is imprimitive in the sense fixed above, i.e.\ a root $\alpha$ of this equation can be regarded as a root of an equation of degree $r$, whose coefficients depend on the root $y$ of an equation of degree $s$.

To prove that then the field $\Omega(\alpha)$ (in the sense of \S.\,144) is also imprimitive, we have still to show that there are quantities in $\Omega(\alpha)$ which satisfy an irreducible equation of degree $s$.

We shall show that $y$ itself is such a quantity, and hence that $y$ can be represented rationally by $\alpha$ alone.

But this is very simple to conclude. By (11) we have
\weq{12}{\psi(\alpha, y) = 0,}
while $\psi(\alpha, y_1), \ldots \psi(\alpha, y_{s-1})$ are different from zero. For if, say, $\psi(\alpha, y_1) = 0$, then $\alpha$ would be a root of the equation of degree $r$, $\psi(t, y_1) = 0$, whose roots are precisely the quantities $\beta, \beta_1 \ldots \beta_{r-1}$, which are different from $\alpha$. Hence the two equations
\weq{13}{\psi(\alpha, u) = 0, \quad \varphi(u) = 0}
have one and only one common root $u = y$. The greatest common divisor of these two functions is therefore linear with respect to $u$, and it gives $y$ rationally in terms of the coefficients of $\varphi(u)$ and $\psi(\alpha, u)$, and \gesperrt{hence rationally in terms of} $\alpha$. With this the converse of 1.\ is proved:

3. \gesperrt{A primitive field has a primitive group.}

\begin{center}
\rule{3cm}{0.4pt}
\end{center}
